\documentclass[12pt]{article}

\usepackage{fullpage}
\usepackage{multicol,multirow}
\usepackage{tabularx}
\usepackage{ulem}
\usepackage[utf8]{inputenc}
\usepackage{mathtools}
\usepackage{pgfplots}
\usepackage[russian]{babel}

\begin{document}

\section*{Лабораторная работа №\,2 по курсу дискрeтного анализа: сбалансированные деревья}

Выполнил студент группы М8О-208Б-22 МАИ \textit{Немкова Анастасия}.

\subsection*{Условие}
Необходимо создать программную библиотеку, реализующую указанную структуру данных, на основе которой разработать программу-словарь. В словаре каждому ключу, представляющему из себя регистронезависимую последовательность букв английского алфавита длиной не более 256 символов, поставлен в соответствие некоторый номер, от 0 до  $2^{64} - 1$. Разным словам может быть поставлен в соответствие один и тот же номер.

Программа должна обрабатывать строки входного файла до его окончания. Каждая строка может иметь следующий формат:

\textbf{+ word 34} — добавить слово «word» с номером 34 в словарь. Программа должна вывести строку «OK», если операция прошла успешно, «Exist», если слово уже находится в словаре.

\textbf{- word} — удалить слово «word» из словаря. Программа должна вывести «OK», если слово существовало и было удалено, «NoSuchWord», если слово в словаре не было найдено.

\textbf{word} — найти в словаре слово «word». Программа должна вывести «OK: 34», если слово было найдено; число, которое следует за «OK:» — номер, присвоенный слову при добавлении. В случае, если слово в словаре не было обнаружено, нужно вывести строку «NoSuchWord».

\textbf{! Save /path/to/file} — сохранить словарь в бинарном компактном представлении на диск в файл, указанный парамером команды. В случае успеха, программа должна вывести «OK», в случае неудачи выполнения операции, программа должна вывести описание ошибки (см. ниже).

\textbf{! Load /path/to/file} — загрузить словарь из файла. Предполагается, что файл был ранее подготовлен при помощи команды Save. В случае успеха, программа должна вывести строку «OK», а загруженный словарь должен заменить текущий (с которым происходит работа); в случае неуспеха, должна быть выведена диагностика, а рабочий словарь должен остаться без изменений. Кроме системных ошибок, программа должна корректно обрабатывать случаи несовпадения формата указанного файла и представления данных словаря во внешнем файле.
\textbf{Задача:}
Требуется разработать программу, осуществляющую ввод пар «ключ-значение», их упорядочивание по возрастанию ключа указанным алгоритмом сортировки за линейное время и вывод отсортированной последовательности.\newline
\newline
\textbf{Вариант:}

    1. AVL-дерево


\subsection*{Метод решения}

Для реализации словаря используется AVL-дерево.\newline
AVL-дерево - это бинарное дерево поиска, являющееся сбалансированным в следующем смысле: высота правого поддерева для любого узла дерева отличается от высоты его левого поддерева не более чем на единицу.\newline
При разнице высот поддеревьев равной 2 или -2 используется балансировка дерева, для этого применяются две вспомогательные функции: правый и левый повороты дерева.\newline
Вставка и удаление узлов происходит как и в обычном бинарном дереве, но в конце процедур выполняется балансировка.\newline
Поиск в AVL-дереве полностью аналогичен поиску в бинарном дереве.\newline
Для записи дерева в файл последовательно записывается размер ключа узла, сам ключ и значение узла, а также информацию о наличии левого и правого потомков.\newline
При загрузке из файла происходит обратный процесс: считывается размер ключа, выделяется память для ключа, считывается сам ключ, затем значение узла, флаги наличия потомков, и, если потомок присутствует, рекурсивно загружается поддерево.\newline
Словарь хранит пары клю-значение. Ключом выступает строка до 256 символов, значение - беззнаковое число от $0$ до $2^{64} - 1$


\subsection*{Описание программы}

Для реализации словаря был создан класс AVL-дерева, каждый узел которого содержит информацию о ключе, значении и высоте узла, а также указатель на правого и левого потомка. Тип данных для значения ключа unsigned long long, позволяющий хранить любое значение до  $2^{64} - 1$ символов. Ключ является строкой перменной длины до 256 символов, для её хранения был создан класс TString, содержащий информацию о размере и вместимости строки и массив элементов типа char. В классе AVL-дерева реализованы следующие методы:

\begin{itemize}
    \item unsigned char GetHeight(TNode* n) - возвращает высоту дерева.
    \item int BFactor(TNode* n) - разница между высотами правого и левого поддерева данного узла.
    \item void FixHeight(TNode* n) - обновление высоты данного узла на основе высот его дочерних узлов.
    \item TNode* RotateRight(TNode* n) - правый поворот вокруг данного узла.
    \item TNode* RotateLeft(TNode* n) - левый поворот вокруг данного узла.
    \item TNode* Balance(TNode* n) - балансировка данного узла в соответствии с его высотой и фактором сбалансированности.
    \item TNode* InsertNode(TNode* n, TNode* k) - вставка узла.
    \item TNode* FindMin(TNode* n) - находит узел с минимальным ключом в дереве с корнем в данном узле.
    \item TNode* RemoveMin(TNode* n) - удаление узла с минимальным ключом в дереве с корнем в данном узле.
    \item TNode* RemoveNode(TNode* n, const T& k) - удаление узла.
    \item TNode* FindNode(TNode* n, const T& k) - нахождение узла с заданным ключом.
    \item void Delete(TNode* n) - удаление дерева.
    \item void Save(TNode* n, std::ofstream& of) - сохранение дерева в файл.
    \item TNode* Load(std::ifstream& is) - загрузка дерева из файла.
\end{itemize}

\subsection*{Дневник отладки}

\begin{enumerate}
    \item 19 апр 2024, 16:07:36 OL на 1 тесте
		
    Решение: добавлена проверка для выхода из цикла,считавающего данные из стандартного потока ввода, при достижении конца файла

    \item 19 апр 2024, 16:11:42 WA на 15 тесте

    Решение: изменено удаление узла дерева: вместо перемещения минимального узла дерева на место удаляемого, создаем новый узел на основе минимального и присоединяем к нему левое поддерево

    \item 24 апр 2024, 00:47:55 TL на тесте 15

    Решение: вместо приведения к нижнему регистру ключа при каждой операции сравнения, используем эту операцию сразу же при считывании значения ключа из стандартного потока ввода

\end{enumerate}

\subsection*{Тест производительности}

Для измерения производительсти сравнивается время операций вставки, удаления и поиска реализованного класса AVL-дерева и стандартного контейнера std::map, в основе которого лежит красно-черное дерево, на одинаковых входных данных. Подсчет времени производился с помощью библиотеки chrono, которая позволяет фиксировать время в начале и конце выполнения сортировки. Первый тест состоит из $10^3$ строк, второй из $5 *10^3$ строк, третий из $10^4$, четвертый из $5 * 10^4$ и пятый из $10^5$ строк.

На графиках представлены зависимости времени выполнения вставки, удаления и поиска узла от объёма входных данных. Сложность данных операций в AVL-дереве $O(log(n))$, где n - количество узлов дерева. Как видно из графика, вставка и удаление элемента работают быстрее в std::map, что связано с необходимостью балансировки AVL-дерева каждый раз после совершения данных операции.


\begin{figure}[htbp]
    \centering
    \begin{tikzpicture}
        \begin{axis}[
            xlabel={Объём входных данных ($\times 10^4$)},
            ylabel={Время работы (ms)},
            grid=major,
            xmin=0, xmax=10,
            ymin=0, ymax=100,
            xtick={0, 1,  2, 3, 4, 5, 6, 7, 8,  9, 10},
            ytick={0, 10, 20, 30, 40, 50, 60, 70, 80, 90, 100},
            legend style={at={(0.5,-0.2)},anchor=north},
            legend columns=2,
            width=0.8\textwidth,
            height=0.5\textwidth,
            ]
            \addplot[color=blue,mark=*] coordinates {
                (0,0)
                (0.1, 0.572)
                (0.5, 3.456)
                (1, 6.987)
                (5, 43.03)
                (10, 99.946)
            };
            \addlegendentry{AVL-tree}
            \addplot[color=red,mark=square] coordinates {
                (0,0)
                (0.1, 0.512)
                (0.5, 2.952)
                (1, 5.778)
                (5, 33.34)
                (10, 73.502)
            };
            \addlegendentry{std::map}
        \end{axis}
    \end{tikzpicture}
    \caption{Вставка}
    \label{fig:graph}
\end{figure}

\begin{figure}[htbp]
    \centering
    \begin{tikzpicture}
        \begin{axis}[
            xlabel={Объём входных данных ($\times 10^4$)},
            ylabel={Время работы (ms)},
            grid=major,
            xmin=0, xmax=10,
            ymin=0, ymax=110,
            xtick={0, 1, 2, 3, 4, 5, 6, 7, 8,  9, 10},
            ytick={0, 11, 22, 33, 44, 55, 66, 77, 88, 99, 111},
            legend style={at={(0.5,-0.2)},anchor=north},
            legend columns=2,
            width=0.8\textwidth,
            height=0.5\textwidth,
            ]
            \addplot[color=blue,mark=*] coordinates {
                (0,0)
                (0.1, 0.477)
                (0.5, 2.747)
                (1, 5.874)
                (5, 43.779)
                (10, 107.007)
            };
            \addlegendentry{AVL-tree}
            \addplot[color=red,mark=square] coordinates {
                (0,0)
                (0.1, 0.544)
                (0.5, 2.679)
                (1, 5.839)
                (5, 40.967)
                (10, 94.293)
            };
            \addlegendentry{std::map}
        \end{axis}
    \end{tikzpicture}
    \caption{Удаление}
    \label{fig:graph}
\end{figure}

\begin{figure}[htbp]
    \centering
    \begin{tikzpicture}
        \begin{axis}[
            xlabel={Объём входных данных ($\times 10^4$)},
            ylabel={Время работы (ms)},
            grid=major,
            xmin=0, xmax=10,
            ymin=0, ymax=4,
            xtick={0, 1, 2, 3, 4, 5, 6, 7, 8,  9, 10},
            ytick={0, 0.4, 0.8, 1.2, 1.6, 2, 2.4, 2.8, 3.2, 3.6, 4},
            legend style={at={(0.5,-0.2)},anchor=north},
            legend columns=2,
            width=0.8\textwidth,
            height=0.5\textwidth,
            ]
            \addplot[color=blue,mark=*] coordinates {
                (0,0)
                (0.1, 0.027)
                (0.5, 0.117)
                (1, 0.24)
                (5, 1.219)
                (10, 2.473)
            };
            \addlegendentry{AVL-tree}
            \addplot[color=red,mark=square] coordinates {
                (0,0)
                (0.1, 0.04)
                (0.5, 0.285)
                (1, 0.46)
                (5, 1.867)
                (10, 3.776)
            };
            \addlegendentry{std::map}
        \end{axis}
    \end{tikzpicture}
    \caption{Поиск}
    \label{fig:graph}
\end{figure}

\subsection*{Выводы}

В ходе данной лабораторной работы былы изучены различные виды деревьев и реализован словарь на основе AVL-дерева. Были разобраны операции работы с данным деревом такие, как вставка, удаления и поиск узла, а также сохранение дерева файл в бинарном представлении и загрузка дерева их файла. Такая структура данных отлично подходит для хранения большого объёма данных, так как все операции выполняются до $O(log(n))$.

\end{document}

